\documentclass[aspectratio=169, 11pt]{beamer}

% ── Theme & colors ──────────────────────────────────────────────────
\usetheme{Madrid}
\usecolortheme{whale}
\definecolor{paris1blue}{RGB}{0,61,121}
\definecolor{accent}{RGB}{59,130,246}
\setbeamercolor{structure}{fg=paris1blue}
\setbeamercolor{title}{fg=white, bg=paris1blue}
\setbeamercolor{frametitle}{fg=white, bg=paris1blue}
\setbeamercolor{block title}{bg=paris1blue, fg=white}
\setbeamercolor{block body}{bg=paris1blue!5}
\setbeamertemplate{navigation symbols}{}
\setbeamertemplate{footline}{%
  \leavevmode\hbox{%
    \begin{beamercolorbox}[wd=.33\paperwidth,ht=2.5ex,dp=1ex,center]{author in head/foot}%
      \usebeamerfont{author in head/foot}\insertshortauthor
    \end{beamercolorbox}%
    \begin{beamercolorbox}[wd=.34\paperwidth,ht=2.5ex,dp=1ex,center]{title in head/foot}%
      \usebeamerfont{title in head/foot}\insertshorttitle
    \end{beamercolorbox}%
    \begin{beamercolorbox}[wd=.33\paperwidth,ht=2.5ex,dp=1ex,right]{date in head/foot}%
      \usebeamerfont{date in head/foot}\insertframenumber{} / \inserttotalframenumber\hspace*{2ex}
    \end{beamercolorbox}}%
}

% ── Packages ────────────────────────────────────────────────────────
\usepackage[utf8]{inputenc}
\usepackage[T1]{fontenc}
\usepackage[french]{babel}
\usepackage{graphicx}
\usepackage{booktabs}
\usepackage{amsmath}
\usepackage{hyperref}
\usepackage{listings}
\lstset{basicstyle=\ttfamily\scriptsize, breaklines=true, frame=single, backgroundcolor=\color{gray!5}}

% ── Metadata ────────────────────────────────────────────────────────
\title[Sentiment Analysis]{Analyse de Sentiment sur Textes Financiers}
\subtitle{Benchmarking Classical ML, Deep Learning et Transformers}
\author[Majed \& Vidal]{Houssem \textsc{Majed} \and Victoria \textsc{Vidal}}
\institute[Paris 1]{Universit\'e Paris 1 Panth\'eon-Sorbonne\\Master MoSEF -- Deep Learning / NLP}
\date{Avril 2026}

\begin{document}

% ════════════════════════════════════════════════════════════════════
% PAGE DE GARDE
% ════════════════════════════════════════════════════════════════════
\begin{frame}[plain]
  \titlepage
\end{frame}

\begin{frame}{Plan}
  \tableofcontents
\end{frame}

% ════════════════════════════════════════════════════════════════════
\section{Introduction}
% ════════════════════════════════════════════════════════════════════

\begin{frame}{Contexte et objectifs}
  \begin{block}{Problematique}
    Les marches financiers produisent un volume massif de textes (communiques de presse,
    rapports d'analystes, depeches). Classifier automatiquement le sentiment de ces
    textes permet d'extraire un signal utile pour l'aide a la decision.
  \end{block}
  \vspace{0.3cm}
  \textbf{Objectif du projet :} comparer trois generations d'approches NLP sur une
  tache de classification de sentiment financier en trois classes (positif, neutre, negatif).
  \vspace{0.3cm}
  \begin{itemize}
    \item Classical ML (TF-IDF + classifieurs lineaires)
    \item Deep Learning (BiLSTM)
    \item Transformers (BERT, DistilBERT, FinBERT)
  \end{itemize}
\end{frame}

% ════════════════════════════════════════════════════════════════════
\section{Dataset}
% ════════════════════════════════════════════════════════════════════

\begin{frame}{FinancialPhraseBank}
  \begin{columns}[T]
    \begin{column}{0.55\textwidth}
      \textbf{Source :} Malo et al. (2014), JASIST 65(4).\\[6pt]
      4\,846 phrases extraites de communiques financiers en anglais,
      annotees par 16 experts selon le point de vue d'un investisseur.\\[6pt]
      \textbf{Question d'annotation :}\\
      \emph{``Cette phrase donne-t-elle une impression positive,
      negative ou neutre de l'entreprise mentionnee ?''}\\[6pt]
      Seules les phrases avec \textbf{accord majoritaire} (>50\%) sont incluses.
    \end{column}
    \begin{column}{0.4\textwidth}
      \begin{table}
        \centering\small
        \begin{tabular}{lrr}
          \toprule
          Classe   & Nombre & Part \\
          \midrule
          Neutral  & 2\,879 & 59.4\% \\
          Positive & 1\,363 & 28.1\% \\
          Negative &    604 & 12.5\% \\
          \bottomrule
        \end{tabular}
      \end{table}
      \vspace{0.3cm}
      {\small\textcolor{red}{Desequilibre important :}\\la classe negative ne represente que 12.5\% du jeu.}
    \end{column}
  \end{columns}
\end{frame}

\begin{frame}{Protocole experimental}
  \begin{block}{Split stratifie}
    70\% train / 10\% validation / 20\% test --- seed=42, stratifie par classe.
  \end{block}
  \vspace{0.3cm}
  \textbf{Metrique principale : Macro-F1}
  \begin{itemize}
    \item Moyenne equitable du F1 par classe
    \item La classe \emph{negative} (12.5\%) pese autant que \emph{neutral} (59.4\%)
    \item Un modele predisant toujours ``neutral'' obtiendrait 59\% d'accuracy mais un macro-F1 proche de 0
  \end{itemize}
  \vspace{0.3cm}
  \textbf{Metrique secondaire : Accuracy} --- reportee pour reference, mais trompeuse sur donnees desequilibrees.
\end{frame}

% ════════════════════════════════════════════════════════════════════
\section{Modeles developpes}
% ════════════════════════════════════════════════════════════════════

\begin{frame}{1. Classical ML --- TF-IDF + Classifieurs}
  \textbf{Pipeline :} texte $\rightarrow$ vecteurs TF-IDF creux (unigrammes + bigrammes, 20K features, TF sublineaire) $\rightarrow$ classifieur.
  \vspace{0.3cm}
  \begin{columns}[T]
    \begin{column}{0.48\textwidth}
      \textbf{Logistic Regression}
      \begin{itemize}
        \item Regularisation L2, solver lbfgs
        \item Rapide a entrainer (<1s)
      \end{itemize}
    \end{column}
    \begin{column}{0.48\textwidth}
      \textbf{Naive Bayes}
      \begin{itemize}
        \item MultinomialNB, alpha=0.1
        \item Hypothese d'independance naive
      \end{itemize}
    \end{column}
  \end{columns}
  \vspace{0.5cm}
  \textbf{Avantages :} rapide, interpretable, pas de GPU necessaire.\\
  \textbf{Limites :} ignore l'ordre des mots, le contexte et la semantique.
\end{frame}

\begin{frame}{2. Deep Learning --- Bidirectional LSTM}
  \textbf{Architecture :} Embedding (64d, 15K vocab) $\rightarrow$ BiLSTM 2 couches (128 unites) $\rightarrow$ Dense $\rightarrow$ Softmax.
  \vspace{0.3cm}
  \begin{itemize}
    \item Entrainement : 20 epochs, batch size 32, Adam optimizer
    \item Early stopping sur la loss de validation
    \item Entrainement sur CPU et GPU (~40s)
  \end{itemize}
  \vspace{0.3cm}
  \textbf{Avantages :} capture l'ordre des mots et les dependances locales.\\[6pt]
  \textbf{Limites :} sans pre-entrainement et avec seulement 5K exemples, le modele ne generalise pas sur la classe minoritaire (0\% recall sur \emph{negative}).
\end{frame}

\begin{frame}{3. Transformers --- BERT, DistilBERT, FinBERT}
  \textbf{Principe :} modeles pre-entraines sur de larges corpus, fine-tunes sur notre dataset.
  \vspace{0.3cm}
  \begin{columns}[T]
    \begin{column}{0.32\textwidth}
      \textbf{BERT}\\[3pt]
      {\small 110M params\\Pre-entraine sur BookCorpus + Wikipedia\\Fine-tuning : 5 epochs, lr=2e-5}
    \end{column}
    \begin{column}{0.32\textwidth}
      \textbf{DistilBERT}\\[3pt]
      {\small 66M params\\Version distillee de BERT (40\% plus legere)\\Meme hyperparametres}
    \end{column}
    \begin{column}{0.32\textwidth}
      \textbf{FinBERT}\\[3pt]
      {\small 110M params\\Further pre-entraine sur corpus financier\\Comprend : \emph{restructuring}, \emph{write-down}, \emph{impairment}...}
    \end{column}
  \end{columns}
  \vspace{0.4cm}
  \textbf{Optimisations GPU :} fp16, gradient checkpointing, VRAM cap a 85\%, gradient accumulation.
\end{frame}

% ════════════════════════════════════════════════════════════════════
\section{Implementation}
% ════════════════════════════════════════════════════════════════════

\begin{frame}{Architecture du projet}
  \begin{columns}[T]
    \begin{column}{0.45\textwidth}
      \textbf{Structure modulaire :}
      \begin{itemize}
        \item \texttt{src/config.py} : chemins, hyperparametres, registre des modeles
        \item \texttt{src/data.py} : chargement, nettoyage, split stratifie
        \item \texttt{src/metrics.py} : evaluation (accuracy, macro-F1, weighted-F1)
        \item \texttt{src/inference.py} : chargement modele, pipeline HuggingFace, prediction
      \end{itemize}
    \end{column}
    \begin{column}{0.5\textwidth}
      \textbf{Scripts d'entrainement :}
      \begin{itemize}
        \item \texttt{train\_baselines.py} : TF-IDF + LogReg/NB
        \item \texttt{train\_lstm.py} : BiLSTM via TensorFlow
        \item \texttt{train\_transformer.py} : BERT/DistilBERT/FinBERT
        \item \texttt{build\_leaderboard.py} : agregation des resultats
      \end{itemize}
      \vspace{0.3cm}
      \textbf{Environnement :}\\
      Python 3.11, uv, PyTorch 2.11+cu128,\\
      GPU NVIDIA RTX 5070 Laptop (8 Go VRAM)
    \end{column}
  \end{columns}
\end{frame}

% ════════════════════════════════════════════════════════════════════
\section{Resultats}
% ════════════════════════════════════════════════════════════════════

\begin{frame}{Benchmark comparatif}
  \begin{table}
    \centering
    \begin{tabular}{clcccc}
      \toprule
      Rang & Modele & Accuracy & Macro-F1 & Type & Temps \\
      \midrule
      1 & \textbf{FinBERT}    & 88.1\% & \textbf{87.7\%} & Transformer & 82s \\
      2 & DistilBERT           & 86.8\% & 86.0\% & Transformer & 51s \\
      3 & BERT                 & 85.5\% & 83.8\% & Transformer & 111s \\
      4 & TF-IDF + NaiveBayes  & 71.3\% & 62.7\% & Classical   & <1s \\
      5 & TF-IDF + LogReg      & 72.4\% & 60.1\% & Classical   & <1s \\
      6 & BiLSTM               & 63.7\% & 37.3\% & Deep Learning & 41s \\
      \bottomrule
    \end{tabular}
  \end{table}
  \vspace{0.3cm}
  \begin{block}{Observations cles}
    \begin{itemize}
      \item Le pre-entrainement sur corpus financier (FinBERT) donne un avantage net
      \item DistilBERT surpasse BERT complet (moins d'overfitting avec peu de donnees)
      \item Les baselines ont une accuracy correcte mais echouent sur la classe minoritaire
      \item Le BiLSTM ne generalise pas sans pre-entrainement
    \end{itemize}
  \end{block}
\end{frame}

% ════════════════════════════════════════════════════════════════════
\section{Dashboard Streamlit}
% ════════════════════════════════════════════════════════════════════

\begin{frame}{Fonctionnalites du dashboard}
  \textbf{Interface web construite avec Streamlit}, deployee localement.
  \vspace{0.3cm}
  \begin{columns}[T]
    \begin{column}{0.48\textwidth}
      \textbf{Analyze}
      \begin{itemize}
        \item Prediction avec un modele au choix
        \item Comparaison multi-modeles cote a cote
        \item Explicabilite LIME (mots cles)
      \end{itemize}
      \vspace{0.2cm}
      \textbf{Batch Analysis}
      \begin{itemize}
        \item Upload CSV de headlines
        \item Classification en masse
        \item Export des resultats
      \end{itemize}
    \end{column}
    \begin{column}{0.48\textwidth}
      \textbf{Live News (Yahoo Finance)}
      \begin{itemize}
        \item API \texttt{yfinance} pour les headlines
        \item Analyse de sentiment en temps reel
        \item Jauge de sentiment global
      \end{itemize}
      \vspace{0.2cm}
      \textbf{Entree vocale}
      \begin{itemize}
        \item Web Speech API (Chrome/Edge)
        \item Injection JS dans le contexte parent
        \item Transcription live dans le textarea
      \end{itemize}
    \end{column}
  \end{columns}
\end{frame}

\begin{frame}{API Yahoo Finance et entree vocale}
  \begin{block}{Integration yfinance}
    On utilise la librairie \texttt{yfinance} pour recuperer les titres d'actualite
    associes a un ticker boursier (AAPL, TSLA, MSFT...). Chaque headline est ensuite
    passee dans le pipeline de prediction du modele actif. On calcule un score de
    sentiment global (proportion positifs $-$ negatifs) pour fournir un indicateur
    synthetique.
  \end{block}
  \vspace{0.3cm}
  \begin{block}{Web Speech API}
    L'entree vocale utilise la Web Speech API du navigateur. Le defi technique principal
    etait que Streamlit execute les composants HTML dans des iframes sandboxes, sans
    acces au micro. La solution retenue injecte un script dans le document parent via
    \texttt{components.html()}, ce qui permet au \texttt{SpeechRecognition} de s'executer
    dans le bon contexte avec les permissions navigateur.
  \end{block}
\end{frame}

% ════════════════════════════════════════════════════════════════════
\section{Difficultes rencontrees}
% ════════════════════════════════════════════════════════════════════

\begin{frame}{Difficultes techniques}
  \begin{enumerate}
    \item \textbf{Gestion de la VRAM} --- Le fine-tuning de BERT (110M params) sur une
      RTX 5070 (8 Go) a necessite fp16, gradient checkpointing et un cap memoire a 85\%.
      Sans ces optimisations, les modeles depassaient la memoire GPU.
    \item \textbf{PyTorch CUDA sur Windows} --- L'installation via \texttt{pip install torch}
      donne une version CPU-only. Il a fallu configurer un index PyTorch specifique
      (\texttt{pytorch-cu128}) dans \texttt{pyproject.toml}.
    \item \textbf{Entree vocale dans Streamlit} --- Les iframes Streamlit bloquent l'acces
      micro. Plusieurs approches ont echoue (\texttt{st\_javascript}, \texttt{components.html}
      directement) avant de trouver la solution d'injection dans le contexte parent.
    \item \textbf{Synchronisation etat Streamlit/DOM} --- Les modifications du textarea par
      JavaScript ne mettaient pas a jour l'etat interne de Streamlit. Resolution via le
      setter React natif et les callbacks \texttt{on\_click}.
    \item \textbf{Desequilibre des classes} --- Le BiLSTM predisait systematiquement
      ``neutral'', obtenant 0\% recall sur \emph{negative}. Ce comportement a confirme
      l'importance du pre-entrainement.
  \end{enumerate}
\end{frame}

% ════════════════════════════════════════════════════════════════════
\section{Conclusion}
% ════════════════════════════════════════════════════════════════════

\begin{frame}{Conclusion et perspectives}
  \begin{block}{Bilan}
    FinBERT domine le benchmark (87.7\% macro-F1) grace a son pre-entrainement
    sur corpus financier. L'ecart avec les baselines classiques confirme que le
    transfert de connaissances est decisif sur un petit jeu de donnees specialise.
  \end{block}
  \vspace{0.3cm}
  \textbf{Perspectives :}
  \begin{itemize}
    \item Elargir le dataset (SEC filings, earnings calls)
    \item Tester des modeles plus recents (DeBERTa, LLMs few-shot)
    \item Deploiement cloud avec API REST
    \item Analyse multi-langue (francais, allemand)
  \end{itemize}
  \vspace{0.5cm}
  \centering\large
  \textbf{Merci de votre attention.}
\end{frame}

\end{document}
